\documentclass{article}
\usepackage{pathfinder-note}

\title{Contracts as a Test Bed for Alignment Mechanisms}

\begin{document}
\maketitle

\section{The pair}

The first paper studies mechanism design for AI agents whose preferences and
capabilities are unknown.  Its abstract distinguishes incentives for honest
reporting from incentives for obedient action, assumes a one-sided capability
structure in which ability can be concealed but not counterfeited, and lists
sandbagging and peer scoring among its applications.  The second paper studies
self-negotiated contracts in a spatial--temporal bargaining and navigation
game.  Its abstract treats observable, enforced agreements as commitment
devices and compares representations ranging from formal contracts that
compile to code to natural contracts that must be reinterpreted
\ledger{7, 11}.

\section{Connexion pursued}

The supported connexion is an experimental one, not a claim that the second
paper already implements the first paper's theory.  Put a controlled
sandbagging problem inside the Contracting Game: assign agents independently
known capability levels with nested feasible action sets, require a capability
report before agreement, and then require a contracted action.  Compare a
no-contract baseline with semantically matched natural-language and formally
executable contracts.  Measure report truthfulness and completion of the
contracted action separately, corresponding respectively to honesty and
obedience \ledger{12, 15}.

The informative result is the boundary between representation and enforcement,
not merely a generic benefit from contracting: do natural contracts match
executable ones, or does reinterpretation permit more concealment or later
disobedience?  This uses the first paper's one-sided imitation structure and
the second paper's explicit formal-to-natural contract axis
\ledger{8, 12, 15}.

\section{Supported findings and objections}

The peer record converges on this design after retracting claims not supported
by the supplied abstracts.  In particular, the available material does not
establish a proper-scoring-regret measure, a transfer formula, detailed
equilibrium conditions, an implementation API, or observed ground truth for
latent agent types \ledger{7, 11}.  Consequently, truthful reports and greater
compliance would support only the instantiated cross-paper mechanism, and only
after its incentives had been checked against the full theory.  Profitable
sandbagging, defection, or a null result would show failure of that
instantiation rather than refute the first paper's theorem
\ledger{9, 12, 16}.

Peer scoring is not the best first intervention on this evidence.  It appears
only as a stylized application in the first abstract, while the second abstract
directly supplies the more diagnostic contrast between executable and
reinterpreted contracts \ledger{13, 16}.  Likewise, compliance cannot stand in
for truthful elicitation: the report stage and action stage need distinct
outcomes \ledger{3, 12}.

\section{Unresolved issues and next step}

Evidence is thin: the supplied inputs are abstracts, and the peer directories
contain no additional notes beyond the ledger record.  The exact mechanism,
the Contracting Game interface, the construction of nested feasible-action
sets, and the semantic equivalence of natural and executable terms remain
unverified \ledger{13, 16}.  The smallest next step is therefore to obtain the
full papers and game code and write a one-page mapping from each theoretical
assumption and incentive to a concrete report, action, payoff, and enforcement
operation.  Only if that mapping succeeds should the three-arm experiment and
its two co-primary outcomes be preregistered.

\end{document}
